
\section{Anomaly-completed topological holography from first principles}
\label{app:anomaly-completed-topological-holography}

\index{topological holography|textbf}
\index{Koszul duality}
\index{chiral algebra}
\index{factorization algebra}
\index{anomaly completion}
\index{transgression algebra}
\index{genus-Clifford completion}
\index{Yangian}
\index{affine Kac--Moody algebra}
\index{Fock module}

Anomaly completion proceeds through three successive projections of $\Theta^{\mathrm{oc}}$: the transgression algebra, the genus-Clifford completion, and the anomaly cancellation $\kappa_{\mathrm{eff}} = \kappa(\text{matter}) + \kappa(\text{ghost}) = 0$. The spectral Drinfeld class vanishes for all simple Lie algebras by root multiplicity one (Theorem~\ref{thm:complete-strictification}); the anomaly data is a separate, genus-dependent phenomenon.

The principal external inputs are:
chiral algebras as quantum coisson structures~\cite{BD04};
chiral/factorization Koszul duality~\cite{FG12};
derived centers in topological field theory~\cite{costello-gaiotto,CP2020,DNP25};
higher operations in HT perturbation theory~\cite{GKW24};
boundary chiral algebras~\cite{CDG20};
Koszul-dual twisted holography~\cite{Zeng23};
and line operators from dg-shifted Yangians~\cite{DNP25}.

\medskip

\noindent\textbf{What is proved here.}
The theorem-level core of this appendix is the construction of:
\begin{enumerate}
\item a universal \emph{transgression algebra} \(B_{\Theta}\) attached to a dg algebra \(B\) and a closed central degree-\(2\) class \(\Theta\);
\item a corrected one-sided two-term resolution computing both \(R\!\operatorname{Hom}\) and derived relative tensor products in the anomaly-completed theory;
\item a \emph{secondary anomaly} \(u=\eta^{2}\), together with a quadratic refinement on the space of neutralizations;
\item a \emph{genus-Clifford completion} \(\mathfrak{G}_{g}(B_{\Theta})\) which packages higher-genus handle operators and exhibits a sharp dichotomy:
after inverting \(u\), higher genus becomes Morita-trivial; modulo \(u\), it degenerates to an exterior extension by \(H_{1}(\Sigma_{g})\).
\end{enumerate}
These statements are proved here and belong to what the monograph elsewhere calls the \emph{M-level}: strict algebraic models.

\medskip

\noindent\textbf{What remains interpretive or conjectural.}
Whenever we identify \(B\) with a bulk Koszul-dual algebra \(A^{!}\), or identify \(\Theta\) with a genuine physical genus anomaly \(\Theta_{g}\), or interpret \(\mathfrak{G}_{g}(B_{\Theta})\) as the full higher-genus line-side package of a twisted field theory, we have crossed from strict algebra into a physical/topological-holographic programme.
Those identifications are supported by \cite{Zeng23,CDG20,DNP25,costello-gaiotto,CP2020}, but outside the strict models below they should be read with appropriate caution.
We will make the distinction explicit as we go.

\subsection{The first-principles correction}
\label{subsec:tholog-first-principles-correction}

One might adjoin to \(B\) an odd \emph{supercentral} generator \(\eta\) with \(d\eta=\Theta\).
This fails.
If \(\eta\) were odd and supercentral in characteristic \(0\), then
\[
\eta^{2}=-\eta^{2},
\]
hence \(\eta^{2}=0\), and the entire secondary theory would collapse before it begins.
The higher-genus refinement would then have no place to live.

The correct object is a \emph{skew polynomial extension}:
\(\eta\) graded-commutes with \(B\) but remains a genuinely free odd generator with respect to its own powers.
Equivalently, one adjoins an odd variable over \(B\) via the sign automorphism
\[
\sigma:B\longrightarrow B,\qquad \sigma(b)=(-1)^{|b|}b.
\]
This is the smallest repair preserving the universal property while leaving room for a nontrivial square
\[
u:=\eta^{2}.
\]

A second correction is equally important.
Because the extension is a skew/Ore extension rather than an ordinary tensor product with an exterior algebra, the natural exact triangle is \emph{one-sided} and module-theoretic.
The correct replacement for the naive untwisted bimodule resolution is a two-term free resolution of each module, from which the transgression formulas for \(R\!\operatorname{Hom}\) and derived tensor products follow cleanly.

\begin{remark}[Why this construction is inevitable]
\label{rem:tholog-inevitability}
The transgression algebra $B_\Theta$ is not one possible
resolution of the curvature obstruction among several: it is
the \emph{universal} one. Its universal property
(Proposition~\ref{prop:tholog-universal-property}) states that
every dg~algebra map $B_\Theta \to C$ is equivalent to a pair
$(f\colon B \to C,\; h \in C^1)$ with $dh = f(\Theta)$ and
$hf(b) = (-1)^{|b|}f(b)h$. In other words, $B_\Theta$
\emph{is} the moduli problem ``resolve $\Theta$ at chain
level'', incarnated as an algebra. The secondary anomaly
$u = \eta^2$ and the genus-Clifford completion
$\mathfrak G_g(B_\Theta)$ are then forced by the same
universality: $u$ is the inevitable by-product of
$\eta^2 \ne 0$, and the Clifford algebra is the universal
algebra in which $2g$ odd generators square to~$u$. The
entire tower (transgression, neutralisation, secondary
anomaly, genus completion, Morita triviality) is a single
deductive chain from the datum $(B, \Theta)$. No choices are
made; no alternatives exist at the same level of universality.
\end{remark}

\subsection{Conventions}
\label{subsec:tholog-conventions}

Throughout this appendix:
\begin{itemize}
\item all complexes are cohomologically graded;
\item the ground field \(k\) has characteristic \(0\);
\item differentials have degree \(+1\);
\item if \(X\) is a graded object, \(|x|\) denotes the degree of a homogeneous element \(x\in X\);
\item for a dg algebra \(B\), \(Z(B)\) denotes its graded center;
\item if \(M\) is a left dg \(B\)-module, \(\underline{\operatorname{Hom}}_{B}(M,N)\) denotes the dg Hom complex;
\item derived functors are written \(R\!\operatorname{Hom}\) and \(\otimes^{L}\);
\item \(\Sigma_{g}\) denotes a closed connected oriented surface of genus \(g\).
\end{itemize}

\subsection{From the source picture to the strict algebraic problem}
\label{subsec:tholog-source-picture}

Let us compress the surrounding story into one strict question.

\begin{quote}
Given a dg algebra \(B\) modelling bulk, defect, or line-operator data, and given a closed central degree-\(2\) anomaly class
\(\Theta\in Z^{2}(B)\cap Z(B)\),
what is the correct universal algebra governing those modules, bimodules, and tensor products on which the action of \(\Theta\) has been neutralized at chain level?
How does one incorporate genus in a way that is invisible away from the anomaly locus and maximally nontrivial on it?
\end{quote}

The answer will be:
\begin{enumerate}
\item first pass from \(B\) to the transgression algebra \(B_{\Theta}\);
\item then, if one wishes to remember genuine surface data, pass further to the genus-Clifford completion \(\mathfrak{G}_{g}(B_{\Theta})\).
\end{enumerate}

Everything below is designed so that these two steps are universal, exact, functorial, and compatible with the familiar monoidal structures appearing in line-operator theories.

\subsection{The transgression algebra}
\label{subsec:tholog-transgression-algebra}

Let \(B\) be a dg algebra, and let
\[
\Theta\in Z^{2}(B)\cap Z(B)
\]
be a closed central degree-\(2\) element.

\begin{definition}[Transgression algebra]
\label{def:tholog-transgression-algebra}
The \emph{transgression algebra} of \((B,\Theta)\) is the dg algebra
\[
B_{\Theta}
:=
\frac{B * k\langle \eta\rangle}
{\langle \eta b-(-1)^{|b|}b\eta\;\mid\; b\in B \text{ homogeneous}\rangle},
\qquad |\eta|=1,\qquad d\eta=\Theta,
\]
where \(B * k\langle\eta\rangle\) denotes the free product of graded algebras and the differential extends that of \(B\).
\end{definition}

Thus \(B_{\Theta}\) is the skew polynomial/Ore extension of \(B\) by the sign automorphism \(\sigma(b)=(-1)^{|b|}b\), equipped with a differential transgressing \(\Theta\).

\begin{remark}[Normal form]
\label{rem:tholog-normal-form}
Every element of \(B_{\Theta}\) may be written uniquely as a finite sum
\[
\sum_{n\geq 0} b_{n}\eta^{n},
\qquad b_{n}\in B,
\]
with all powers of \(\eta\) placed on the right.
Equivalently, \(B_{\Theta}\) is free as both a left and a right \(B\)-module on the basis \(\{1,\eta,\eta^{2},\dots\}\).
\end{remark}

\begin{proposition}[Universal property; \ClaimStatusProvedHere]
\label{prop:tholog-universal-property}
Let \(C\) be a dg algebra.
To give a dg algebra map
\[
\widetilde f:B_{\Theta}\longrightarrow C
\]
is equivalent to giving:
\begin{enumerate}
\item a dg algebra map \(f:B\to C\);
\item a degree-\(1\) element \(h\in C^{1}\)
\end{enumerate}
such that
\[
h\,f(b)=(-1)^{|b|}f(b)\,h
\qquad \text{for all homogeneous } b\in B,
\]
and
\[
dh=f(\Theta).
\]
Under this correspondence one has \(\widetilde f(\eta)=h\).
\end{proposition}

\begin{proof}
A dg algebra map \(\widetilde f\) is determined by its restriction \(f=\widetilde f|_{B}\) and by the image
\(h=\widetilde f(\eta)\).
Preservation of the defining relations for \(B_{\Theta}\) is exactly the relation
\[
h\,f(b)=(-1)^{|b|}f(b)\,h,
\]
and compatibility with the differential is exactly
\[
dh=d\widetilde f(\eta)=\widetilde f(d\eta)=\widetilde f(\Theta)=f(\Theta).
\]
Conversely, any pair \((f,h)\) satisfying those two conditions extends uniquely to a dg algebra map on the free product and descends through the quotient by the relation \(\eta b=(-1)^{|b|}b\eta\).
\end{proof}

\begin{corollary}[Modules as neutralizations; \ClaimStatusProvedHere]
\label{cor:tholog-modules-neutralizations}
To give a left dg \(B_{\Theta}\)-module structure on a graded vector space \(M\) is equivalent to giving:
\begin{enumerate}
\item a left dg \(B\)-module structure on \(M\);
\item a degree-\(1\) \(B\)-linear endomorphism
\[
h_{M}\in \underline{\operatorname{End}}^{1}_{B}(M)
\]
such that
\[
d_{\underline{\operatorname{End}}}(h_{M})=\Theta_{M},
\]
where \(\Theta_{M}\) denotes the action of \(\Theta\) on \(M\).
\end{enumerate}
\end{corollary}

\begin{proof}
Apply Proposition~\ref{prop:tholog-universal-property} to
\[
C=\underline{\operatorname{End}}(M)
\]
and note that the \(B\)-centrality condition becomes precisely \(B\)-linearity of the degree-\(1\) endomorphism \(h_{M}\).
\end{proof}

\begin{definition}[Neutralizable module]
\label{def:tholog-neutralizable}
A left dg \(B\)-module \(M\) will be called \(\Theta\)-\emph{neutralizable} if the action of \(\Theta\) on \(M\) is null-homotopic, i.e.\ if there exists
\[
h_{M}\in \underline{\operatorname{End}}^{1}_{B}(M)
\]
with
\[
d_{\underline{\operatorname{End}}}(h_{M})=\Theta_{M}.
\]
Any such \(h_{M}\) will be called a \emph{neutralization} of the \(\Theta\)-action on \(M\).
\end{definition}

\begin{corollary}[Obstruction class; \ClaimStatusProvedHere]
\label{cor:tholog-obstruction-class}
A left dg \(B\)-module \(M\) is \(\Theta\)-neutralizable if and only if
\[
[\Theta_{M}]=0\in H^{2}\underline{\operatorname{End}}_{B}(M)
=\operatorname{Ext}^{2}_{B}(M,M).
\]
\end{corollary}

\begin{proof}
By definition, \(M\) is neutralizable if and only if \(\Theta_{M}\) is a coboundary in the dg algebra \(\underline{\operatorname{End}}_{B}(M)\).
\end{proof}

\begin{corollary}[Moduli of neutralizations; \ClaimStatusProvedHere]
\label{cor:tholog-moduli-neutralizations}
If \(M\) is \(\Theta\)-neutralizable, then the set of homotopy classes of neutralizations of \(M\) is an affine space over
\[
H^{1}\underline{\operatorname{End}}_{B}(M)
=\operatorname{Ext}^{1}_{B}(M,M).
\]
\end{corollary}

\begin{proof}
If \(h_{M}\) and \(h_{M}'\) are two neutralizations, then
\[
d(h_{M}'-h_{M})=0.
\]
Thus the difference is a degree-\(1\) cocycle.
Changing a neutralization by a boundary does not change its homotopy class.
\end{proof}

\subsection{The corrected one-sided resolution}
\label{subsec:tholog-corrected-one-sided-resolution}

The next theorem replaces the naive untwisted bimodule resolution.
It is the technical engine behind the exact formulas for derived morphisms and relative tensor products.

\begin{theorem}[Free two-term resolution of a neutralized module; \ClaimStatusProvedHere]
\label{thm:tholog-free-two-term-resolution}
Let \(M\) be a left dg \(B_{\Theta}\)-module, and let \(h_{M}\) be the corresponding neutralization of the underlying left dg \(B\)-module.
Define a degree-\(1\) map of graded left \(B_{\Theta}\)-modules
\[
D_{M}:B_{\Theta}\otimes_{B} M\longrightarrow B_{\Theta}\otimes_{B} M
\]
by
\[
D_{M}(s\otimes m)=s\eta\otimes m-(-1)^{|s|}s\otimes h_{M}(m).
\]
Then \(D_{M}\) is a chain map of degree \(1\), and there is a short exact sequence of dg \(B_{\Theta}\)-modules
\[
0\longrightarrow \big(B_{\Theta}\otimes_{B}M\big)[-1]
\xrightarrow{\;\widetilde D_{M}\;}
B_{\Theta}\otimes_{B}M
\xrightarrow{\;\mu_{M}\;}
M
\longrightarrow 0,
\]
where \(\widetilde D_{M}\) is the corresponding degree-\(0\) map from the shifted domain and \(\mu_{M}(s\otimes m)=s\cdot m\).
\end{theorem}

\begin{proof}
We first check that \(D_{M}\) is well defined over \(B\).
For homogeneous \(b\in B\),
\[
D_{M}(sb\otimes m)
=
sb\eta\otimes m-(-1)^{|s|+|b|}sb\otimes h_{M}(m).
\]
Since \(b\eta=(-1)^{|b|}\eta b\) in \(B_{\Theta}\),
\[
sb\eta\otimes m=(-1)^{|b|}s\eta b\otimes m=s\eta\otimes bm,
\]
and because \(h_{M}\) is \(B\)-linear,
\[
sb\otimes h_{M}(m)=s\otimes bh_{M}(m)=s\otimes h_{M}(bm).
\]
Hence
\[
D_{M}(sb\otimes m)=D_{M}(s\otimes bm).
\]

We next show that \(D_{M}\) is a chain map.
Write \(d\) for the differential on \(B_{\Theta}\otimes_{B}M\).
A direct computation gives
\begin{align*}
d\,D_{M}(s\otimes m)
&=
d(s\eta)\otimes m+(-1)^{|s|+1}s\eta\otimes dm
\\
&\qquad
-(-1)^{|s|}d(s)\otimes h_{M}(m)-s\otimes d(h_{M}m),
\end{align*}
while
\begin{align*}
D_{M}\,d(s\otimes m)
&=
D_{M}\big(d(s)\otimes m+(-1)^{|s|}s\otimes dm\big)
\\
&=
d(s)\eta\otimes m-(-1)^{|s|+1}d(s)\otimes h_{M}(m)
\\
&\qquad
+(-1)^{|s|}s\eta\otimes dm-s\otimes h_{M}(dm).
\end{align*}
Subtracting, using \(d(s\eta)=d(s)\eta+(-1)^{|s|}s\Theta\) and
\[
d(h_{M}m)=d_{\underline{\operatorname{End}}}(h_{M})(m)-h_{M}(dm)=\Theta_{M}(m)-h_{M}(dm),
\]
we obtain
\[
dD_{M}(s\otimes m)+D_{M}d(s\otimes m)=
(-1)^{|s|}s\Theta\otimes m-s\otimes \Theta_{M}(m)=0.
\]
So \(D_{M}\) is a chain map.

\(\mu_{M}\) is surjective: \(\mu_M(1 \otimes m) = m\).
The image of \(D_{M}\) lies in \(\ker(\mu_{M})\) because
\[
\mu_{M}\big(D_{M}(s\otimes m)\big)
=s\eta\cdot m-(-1)^{|s|}s\cdot h_{M}(m)=0.
\]
Conversely, \(B_{\Theta}\otimes_{B}M\) is the free left \(B_{\Theta}\)-module on the underlying \(B\)-module \(M\), and \(M\) is obtained from it by imposing the relations
\[
\eta\otimes m=1\otimes h_{M}(m).
\]
Hence \(\ker(\mu_{M})\) is generated by the image of \(D_{M}\).

It remains to prove injectivity of \(\widetilde D_{M}\), equivalently of \(D_{M}\) as a degree-\(1\) map on the underlying graded modules.
Write an element in normal form:
\[
x=\sum_{i=0}^{N}\eta^{i}\otimes m_{i}.
\]
Then
\[
D_{M}(x)=\sum_{i=0}^{N}\eta^{i+1}\otimes m_{i}-\sum_{i=0}^{N}(-1)^{i}\eta^{i}\otimes h_{M}(m_{i}).
\]
The coefficient of \(\eta^{N+1}\) is \(m_{N}\), so if \(D_{M}(x)=0\), then \(m_{N}=0\).
Descending induction on \(N\) gives \(m_{i}=0\) for all \(i\).
Thus \(D_{M}\) is injective.
\end{proof}

\begin{remark}[Why one-sided is the right level]
\label{rem:tholog-why-one-sided}
Theorem~\ref{thm:tholog-free-two-term-resolution} is the precise place where the Ore nature of \(B_{\Theta}\) matters.
One should think of \(B_{\Theta}\) not as \(B\) tensored with an exterior algebra, but as a skew polynomial extension.
The resulting exact triangle is naturally attached to a \emph{module}, not first to a bimodule.
This is the correct algebraic setting for the transgression formulas below.
\end{remark}

\subsection{Transgression of derived morphisms}
\label{subsec:tholog-derived-morphisms}

\begin{theorem}[Transgression exact triangle for derived Hom; \ClaimStatusProvedHere]
\label{thm:tholog-derived-hom}
Let \(M,N\) be left dg \(B_{\Theta}\)-modules, with neutralizations \(h_{M}\) and \(h_{N}\).
Define the degree-\(1\) endomorphism
\[
\delta_{M,N}:\underline{\operatorname{Hom}}_{B}(M,N)\longrightarrow
\underline{\operatorname{Hom}}_{B}(M,N)
\]
by
\[
\delta_{M,N}(f)=h_{N}\circ f-(-1)^{|f|}f\circ h_{M}.
\]
Then there is a functorial exact triangle
\[
R\!\operatorname{Hom}_{B_{\Theta}}(M,N)
\longrightarrow
R\!\operatorname{Hom}_{B}(M,N)
\xrightarrow{\;\delta_{M,N}\;}
R\!\operatorname{Hom}_{B}(M,N)[1]
\longrightarrow .
\]
Equivalently,
\[
R\!\operatorname{Hom}_{B_{\Theta}}(M,N)
\simeq
\operatorname{Fib}\!\left(
R\!\operatorname{Hom}_{B}(M,N)
\xrightarrow{\;\delta_{M,N}\;}
R\!\operatorname{Hom}_{B}(M,N)[1]
\right).
\]
\end{theorem}

\begin{proof}
Apply \(\underline{\operatorname{Hom}}_{B_{\Theta}}(-,N)\) to the short exact sequence of Theorem~\ref{thm:tholog-free-two-term-resolution}.
Because \(B_{\Theta}\otimes_{B}M\) is a free left \(B_{\Theta}\)-module on the underlying \(B\)-module \(M\), the standard adjunction gives
\[
\underline{\operatorname{Hom}}_{B_{\Theta}}(B_{\Theta}\otimes_{B}M,N)
\cong
\underline{\operatorname{Hom}}_{B}(M,N).
\]
Under this identification, precomposition by \(D_{M}\) becomes exactly
\[
f\longmapsto h_{N}f-(-1)^{|f|}fh_{M}.
\]
Passing to derived categories yields the exact triangle.
\end{proof}

\begin{corollary}[Long exact sequence in Ext; \ClaimStatusProvedHere]
\label{cor:tholog-long-exact-ext}
For any left dg \(B_{\Theta}\)-modules \(M,N\), there is a natural long exact sequence
\[
\cdots\to
\operatorname{Ext}^{n}_{B_{\Theta}}(M,N)
\to
\operatorname{Ext}^{n}_{B}(M,N)
\xrightarrow{\;\delta_{*}\;}
\operatorname{Ext}^{n+1}_{B}(M,N)
\to
\operatorname{Ext}^{n+1}_{B_{\Theta}}(M,N)
\to\cdots.
\]
\end{corollary}

\begin{remark}[The philosophy]
\label{rem:tholog-hom-philosophy}
The anomaly-completed morphism complex is computed from the underlying \(B\)-morphism complex by imposing one degree-\(1\) transgression operator \(\delta_{M,N}\).
Anomaly-completed Ext groups are controlled by ordinary Ext groups plus one universal connecting morphism.
\end{remark}

\subsection{Transgression of derived relative tensor products}
\label{subsec:tholog-derived-tensors}

For a right \(B_{\Theta}\)-module, the neutralization operator is not strictly right \(B\)-linear; it is twisted by the sign automorphism.
We therefore record the correct right-handed notion.

\begin{definition}[Right neutralization]
\label{def:tholog-right-neutralization}
A right dg \(B_{\Theta}\)-module \(L\) is equivalently a right dg \(B\)-module together with a degree-\(1\) endomorphism
\[
h_{L}:L\longrightarrow L
\]
such that
\[
h_{L}(lb)=(-1)^{|b|}h_{L}(l)b
\qquad (b\in B \text{ homogeneous}),
\]
and
\[
d(h_{L})=\Theta_{L},
\]
where \(\Theta_{L}\) denotes right multiplication by \(\Theta\).
We call \(h_{L}\) a \emph{right neutralization}.
\end{definition}

\begin{theorem}[Transgression exact triangle for derived tensor products; \ClaimStatusProvedHere]
\label{thm:tholog-derived-tensor}
Let \(L\) be a right dg \(B_{\Theta}\)-module and \(R\) a left dg \(B_{\Theta}\)-module, with right and left neutralizations \(h_{L}\) and \(h_{R}\).
Then there is a functorial exact triangle
\[
L\otimes^{L}_{B}R[-1]
\xrightarrow{\;\theta_{L,R}\;}
L\otimes^{L}_{B}R
\longrightarrow
L\otimes^{L}_{B_{\Theta}}R
\longrightarrow ,
\]
where on homogeneous tensors
\[
\theta_{L,R}(l\otimes r)
=
(-1)^{|l|}h_{L}(l)\otimes r-l\otimes h_{R}(r).
\]
Equivalently,
\[
L\otimes^{L}_{B_{\Theta}}R
\simeq
\operatorname{Cone}\!\left(
L\otimes^{L}_{B}R[-1]
\xrightarrow{\;\theta_{L,R}\;}
L\otimes^{L}_{B}R
\right).
\]
\end{theorem}

\begin{proof}
Resolve \(R\) by the free two-term resolution from Theorem~\ref{thm:tholog-free-two-term-resolution}, then tensor on the left over \(B_{\Theta}\) with \(L\).
The resulting two-term complex is
\[
L\otimes_{B}R[-1]\xrightarrow{\theta_{L,R}}L\otimes_{B}R,
\]
and a direct inspection shows that the induced map is exactly
\[
l\otimes r\longmapsto (-1)^{|l|}h_{L}(l)\otimes r-l\otimes h_{R}(r).
\]
The sign in the first term is forced by the \(\sigma\)-semilinearity of the right neutralization and ensures that the map is \(B\)-balanced.
Passing to derived tensor products yields the triangle.
\end{proof}

\begin{remark}[Interval compactification]
\label{rem:tholog-interval-compactification}
Theorem~\ref{thm:tholog-derived-tensor} is the strict algebraic shadow of the interval/strip compactification formula familiar from boundary topological field theory.
At the level of strict dg models, the first anomaly correction to an interval compactification is always a universal two-term cone.
This is exactly the form one would want when comparing boundary conditions through a bulk or line algebra endowed with a nontrivial central genus class.
\end{remark}

\subsection{Gauge invariance of the transgression operator}
\label{subsec:tholog-gauge-invariance}

Neutralizations are never unique.
The next theorem says that the transgression operator on morphism complexes depends only on the \(B_{\Theta}\)-module structures, not on the chosen representatives.

\begin{theorem}[Gauge invariance of transgression; \ClaimStatusProvedHere]
\label{thm:tholog-gauge-invariance}
Let \(M,N\) be fixed left dg \(B\)-modules, and suppose
\[
h_{M}'=h_{M}+d(u_{M}),
\qquad
h_{N}'=h_{N}+d(u_{N}),
\]
for degree-\(0\) \(B\)-linear endomorphisms \(u_{M},u_{N}\).
Then the corresponding operators
\[
\delta_{M,N},\delta_{M,N}' :
\underline{\operatorname{Hom}}_{B}(M,N)\to
\underline{\operatorname{Hom}}_{B}(M,N)[1]
\]
are chain-homotopic.
\end{theorem}

\begin{proof}
Define a degree-\(0\) endomorphism
\[
H(f)=u_{N}\circ f-f\circ u_{M}.
\]
A direct computation in the dg algebra \(\underline{\operatorname{Hom}}_{B}(M,N)\) gives
\[
\delta_{M,N}'-\delta_{M,N}=dH+Hd.
\]
Hence the two operators are chain-homotopic.
\end{proof}

\subsection{Monoidal lift for primitive anomalies}
\label{subsec:tholog-monoidal-lift}

The next theorem explains when the anomaly-completed theory inherits a monoidal structure.

\begin{definition}[Primitive anomaly]
\label{def:tholog-primitive-anomaly}
Assume that \(B\) is a dg bialgebra with coproduct \(\Delta\) and counit \(\varepsilon\).
A closed central degree-\(2\) class
\[
\Theta\in Z^{2}(B)\cap Z(B)
\]
will be called \emph{primitive} if
\[
\Delta(\Theta)=\Theta\otimes 1+1\otimes \Theta,
\qquad
\varepsilon(\Theta)=0.
\]
\end{definition}

\begin{theorem}[Primitive lift theorem; \ClaimStatusProvedHere]
\label{thm:tholog-primitive-lift}
Let \(B\) be a dg bialgebra and let \(\Theta\in Z^{2}(B)\cap Z(B)\) be primitive.
Then the transgression algebra \(B_{\Theta}\) admits a unique dg bialgebra structure extending that of \(B\) and satisfying
\[
\Delta(\eta)=\eta\otimes 1+1\otimes \eta,
\qquad
\varepsilon(\eta)=0.
\]
\end{theorem}

\begin{proof}
Because \(B_{\Theta}\) is generated as a graded algebra by \(B\) and \(\eta\), uniqueness is immediate.

Define \(\Delta\) and \(\varepsilon\) on generators by the given formulas.
It remains to check that the defining relation
\[
\eta b=(-1)^{|b|}b\eta
\]
is preserved.
Write
\[
\Delta(b)=\sum b_{(1)}\otimes b_{(2)}.
\]
Then in the graded tensor product algebra,
\begin{align*}
\Delta(\eta)\Delta(b)
&=
(\eta\otimes 1+1\otimes \eta)\sum b_{(1)}\otimes b_{(2)}
\\
&=
\sum \eta b_{(1)}\otimes b_{(2)}+\sum (-1)^{|b_{(1)}|}b_{(1)}\otimes \eta b_{(2)}
\\
&=
\sum (-1)^{|b_{(1)}|}b_{(1)}\eta\otimes b_{(2)}
+\sum (-1)^{|b|}b_{(1)}\otimes b_{(2)}\eta
\\
&=
(-1)^{|b|}\Delta(b)\Delta(\eta).
\end{align*}
So the relation is respected.

Compatibility with the differential on \(\eta\) is the primitive condition:
\[
d\,\Delta(\eta)
=
\Theta\otimes 1+1\otimes \Theta
=
\Delta(\Theta)
=
\Delta(d\eta).
\]
All other identities are inherited from \(B\).
\end{proof}

\begin{corollary}[Fusion of neutralizations; \ClaimStatusProvedHere]
\label{cor:tholog-fusion-neutralizations}
Assume the setting of Theorem~\ref{thm:tholog-primitive-lift}.
If \(M,N\) are left dg \(B_{\Theta}\)-modules with neutralizations \(h_{M}\) and \(h_{N}\), then \(M\otimes N\) becomes a left dg \(B_{\Theta}\)-module with neutralization
\[
h_{M\otimes N}=h_{M}\otimes 1+1\otimes h_{N},
\]
where \(1\otimes h_{N}\) is understood in the graded sense, so that on homogeneous tensors
\[
h_{M\otimes N}(m\otimes n)=h_{M}(m)\otimes n+(-1)^{|m|}m\otimes h_{N}(n).
\]
\end{corollary}

\begin{proof}
This is the action induced by \(\Delta(\eta)=\eta\otimes 1+1\otimes \eta\).
Since \(\eta\) is primitive, the differential condition follows from \(d(h_{M})=\Theta_{M}\) and \(d(h_{N})=\Theta_{N}\).
\end{proof}

\subsection{The secondary anomaly}
\label{subsec:tholog-secondary-anomaly}

The central square of the transgression generator is invisible in naive supercentral models and inevitable in the corrected skew model.

\begin{proposition}[The secondary central class; \ClaimStatusProvedHere]
\label{prop:tholog-secondary-central-class}
Let
\[
u:=\eta^{2}\in B_{\Theta}^{2}.
\]
Then \(u\) is a closed central degree-\(2\) element of \(B_{\Theta}\).
\end{proposition}

\begin{proof}
Since \(d\eta=\Theta\) and \(\Theta\) has even degree,
\[
d(u)=d(\eta^{2})=(d\eta)\eta-\eta(d\eta)=\Theta\eta-\eta\Theta=0,
\]
because \(\eta\) graded-commutes with \(B\), hence with \(\Theta\).

To check centrality, let \(b\in B\) be homogeneous.
Then
\[
u\,b=\eta(\eta b)=(-1)^{|b|}\eta(b\eta)=(-1)^{2|b|}b\eta^{2}=bu.
\]
Since \(u = \eta^2\), associativity gives \(u\eta = \eta^3 = \eta u\), so \(u\) commutes with~\(\eta\) and hence with all of~\(B_\Theta\).
\end{proof}

\begin{definition}[Secondary anomaly of a lift]
\label{def:tholog-secondary-anomaly}
Let \(M\) be a \(\Theta\)-neutralizable left dg \(B\)-module and let \(h\) be a chosen neutralization.
The \emph{secondary anomaly} of the lift \((M,h)\) is the cohomology class
\[
\sigma(M,h):=[h^{2}]
\in
H^{2}\underline{\operatorname{End}}_{B}(M)
=
\operatorname{Ext}^{2}_{B}(M,M).
\]
Equivalently, \(\sigma(M,h)\) is the action of the central element \(u=\eta^{2}\) on \(M\).
\end{definition}

\begin{remark}[Why this is genuinely new]
\label{rem:tholog-secondary-why-new}
The secondary anomaly vanishes if one forces \(\eta\) to be odd supercentral.
In the corrected skew extension, \(u=\eta^{2}\) is closed, central, and nontrivial:
the place where higher genus lands.
\end{remark}

\subsection{Quadratic geometry of the space of neutralizations}
\label{subsec:tholog-quadratic-geometry}

The moduli of neutralizations carries a canonical quadratic refinement valued in degree-\(2\) self-Ext.

\begin{theorem}[Quadratic refinement on the space of neutralizations; \ClaimStatusProvedHere]
\label{thm:tholog-quadratic-refinement}
Let \(M\) be a \(\Theta\)-neutralizable left dg \(B\)-module, and choose one neutralization \(h\).
Let
\[
\mathcal{L}_{\Theta}(M)
:=
\big\{
h'\in \underline{\operatorname{End}}^{1}_{B}(M)\;\big|\;
d(h')=\Theta_{M}
\big\}.
\]
Then:
\begin{enumerate}
\item \(\mathcal{L}_{\Theta}(M)\) is an affine space over the degree-\(1\) cocycles
\[
Z^{1}\underline{\operatorname{End}}_{B}(M).
\]
\item The assignment
\[
q_{h}(z):=[(h+z)^{2}-h^{2}]
\]
descends to a well-defined map
\[
q_{h}:H^{1}\underline{\operatorname{End}}_{B}(M)\longrightarrow
H^{2}\underline{\operatorname{End}}_{B}(M),
\]
and admits the explicit formula
\[
q_{h}([z])=[hz+zh+z^{2}].
\]
\item Its polarization
\[
b([z],[w])
=
q_{h}([z]+[w])-q_{h}([z])-q_{h}([w])
\]
is independent of the chosen base point \(h\) and is given by
\[
b([z],[w])=[zw+wz].
\]
Thus \(H^{1}\underline{\operatorname{End}}_{B}(M)\) carries a canonical symmetric bilinear form with values in \(H^{2}\underline{\operatorname{End}}_{B}(M)\), and each neutralization \(h\) provides a quadratic refinement of that form.
\end{enumerate}
\end{theorem}

\begin{proof}
If \(h',h''\in \mathcal{L}_{\Theta}(M)\), then
\[
d(h''-h')=0,
\]
so \(\mathcal{L}_{\Theta}(M)\) is an affine space over \(Z^{1}\underline{\operatorname{End}}_{B}(M)\).
This proves (1).

For (2), if \(z\in Z^{1}\underline{\operatorname{End}}_{B}(M)\), then \(h+z\) is again a neutralization, so \((h+z)^{2}-h^{2}=hz+zh+z^{2}\) is a degree-\(2\) cocycle.
We must show that its cohomology class depends only on \([z]\).
Replace \(z\) by \(z+d(a)\) with \(a\in \underline{\operatorname{End}}^{0}_{B}(M)\).
Set \(x=h+z\), so \(dx=\Theta_{M}\).
Then
\[
(x+d(a))^{2}-x^{2}=x\,d(a)+d(a)\,x+d(a)^{2}.
\]
Since \(d^{2}(a)=[\Theta_{M},a]=0\) (because \(\Theta_{M}\) comes from the action of the central element \(\Theta\in B\)), one computes
\[
x\,d(a)+d(a)\,x+d(a)^{2}
=
d\big(-xa+ax+a\,d(a)\big).
\]
Thus the cohomology class is unchanged.
This proves (2).

Finally, polarization gives
\[
q_{h}(z+w)-q_{h}(z)-q_{h}(w)=zw+wz,
\]
which is independent of \(h\).
This proves (3).
\end{proof}

\subsection{Curvature identity on morphism complexes}
\label{subsec:tholog-curvature-identity}

The transgression operator \(\delta_{M,N}\) is rarely square-zero; the secondary anomaly measures the failure exactly.

\begin{theorem}[Curvature identity; \ClaimStatusProvedHere]
\label{thm:tholog-curvature-identity}
Let \(M,N\) be left dg \(B_{\Theta}\)-modules.
Then the operator
\[
\delta_{M,N}(f)=h_{N}f-(-1)^{|f|}fh_{M}
\]
from Theorem~\ref{thm:tholog-derived-hom} is a degree-\(1\) chain map satisfying
\[
\delta_{M,N}^{2}(f)=h_{N}^{2}f-fh_{M}^{2}.
\]
Equivalently,
\[
\delta_{M,N}^{2}
=
u_{N}\circ(-)-(-)\circ u_{M},
\]
where \(u_{M}=h_{M}^{2}\) and \(u_{N}=h_{N}^{2}\).
Passing to cohomology,
\[
[\delta_{M,N}]^{2}
=
\sigma(N,h_{N})\circ(-)-(-)\circ \sigma(M,h_{M})
\]
on \(\operatorname{Ext}^{\bullet}_{B}(M,N)\).
\end{theorem}

\begin{proof}
A direct computation shows that \(\delta_{M,N}\) is a chain map because
\[
d(h_{N})=\Theta_{N},\qquad d(h_{M})=\Theta_{M},
\]
and the \(\Theta\)-terms cancel by centrality.
Squaring \(\delta_{M,N}\), we obtain
\begin{align*}
\delta_{M,N}^{2}(f)
&=
h_{N}\big(h_{N}f-(-1)^{|f|}fh_{M}\big)
-
(-1)^{|f|+1}\big(h_{N}f-(-1)^{|f|}fh_{M}\big)h_{M}
\\
&=
h_{N}^{2}f-fh_{M}^{2},
\end{align*}
because the middle terms cancel.
The cohomological statement follows.
\end{proof}

\begin{corollary}[Orthogonality of sectors with distinct secondary anomaly; \ClaimStatusProvedHere]
\label{cor:tholog-orthogonality}
Let \(M,N\) be left dg \(B_{\Theta}\)-modules.
If the degree-\(2\) endomorphism
\[
u_{N}\circ(-)-(-)\circ u_{M}
\]
acts by a homotopy-invertible endomorphism on \(R\!\operatorname{Hom}_{B}(M,N)\), then
\[
R\!\operatorname{Hom}_{B_{\Theta}}(M,N)\simeq 0.
\]
\end{corollary}

\begin{proof}
By Theorem~\ref{thm:tholog-curvature-identity}, the transgression operator \(\delta_{M,N}\) has square equal to the displayed endomorphism.
If that square is homotopy invertible, then \(\delta_{M,N}\) itself is homotopy invertible, with homotopy inverse
\[
\big(u_{N}\circ(-)-(-)\circ u_{M}\big)^{-1}\delta_{M,N}.
\]
But the fiber of a homotopy equivalence is acyclic, so Theorem~\ref{thm:tholog-derived-hom} gives the claim.
\end{proof}

\begin{corollary}[Equal-secondary sector; \ClaimStatusProvedHere]
\label{cor:tholog-equal-secondary-sector}
If \(u_{M}=u_{N}\) as chain maps, then
\[
\delta_{M,N}^{2}=0.
\]
Hence \(\underline{\operatorname{Hom}}_{B}(M,N)\) becomes a mixed complex with differentials \(d\) and \(\delta_{M,N}\), and one has a short exact sequence
\[
0\longrightarrow
\operatorname{coker}\!\big(H^{n-1}\delta_{M,N}\big)
\longrightarrow
\operatorname{Ext}^{n}_{B_{\Theta}}(M,N)
\longrightarrow
\ker\!\big(H^{n}\delta_{M,N}\big)
\longrightarrow 0.
\]
\end{corollary}

\subsection{Cohomology of the transgression algebra as a hypersurface}
\label{subsec:tholog-hypersurface}

The transgression algebra admits a very concrete cohomological description.

Let
\[
C:=B[u],\qquad |u|=2,\qquad du=0.
\]
As a graded left \(C\)-module, every element of \(B_{\Theta}\) has a unique decomposition
\[
c_{0}+c_{1}\eta,
\qquad c_{0},c_{1}\in C.
\]

\begin{theorem}[Hypersurface long exact sequence; \ClaimStatusProvedHere]
\label{thm:tholog-hypersurface}
There is a short exact sequence of complexes of left \(C\)-modules
\[
0\longrightarrow C
\longrightarrow B_{\Theta}
\longrightarrow C[-1]
\longrightarrow 0,
\]
whose connecting morphism on cohomology is multiplication by \([\Theta]\) (up to the standard cohomological sign convention).
Consequently, there is a long exact sequence
\[
\cdots\to
H^{n}(C)\to H^{n}(B_{\Theta})\to H^{n-1}(C)
\xrightarrow{\cdot[\Theta]}
H^{n+1}(C)\to\cdots.
\]
If multiplication by \([\Theta]\) is injective on \(H^{\bullet}(B)[u]\), then
\[
H^{\bullet}(B_{\Theta})\cong H^{\bullet}(B)[u]/([\Theta]).
\]
\end{theorem}

\begin{proof}
The inclusion \(C\hookrightarrow B_{\Theta}\) is the inclusion of the even part \(\sum b_{n}u^{n}\), and the quotient identifies with \(C[-1]\) via
\[
c\longmapsto c\eta.
\]
So there is a short exact sequence of graded \(C\)-modules.
The differential on \(c\eta\) is
\[
d(c\eta)=d(c)\eta+(-1)^{|c|}c\Theta.
\]
Hence the connecting morphism sends the class of \(c\) to the class of \(c\Theta\), up to the standard sign from the shift.
This yields the long exact sequence.
If multiplication by \([\Theta]\) is injective, exactness identifies \(H^{\bullet}(B_{\Theta})\) with the quotient of \(H^{\bullet}(C)=H^{\bullet}(B)[u]\) by the principal ideal generated by \([\Theta]\).
\end{proof}

\begin{remark}[The hypersurface viewpoint]
\label{rem:tholog-hypersurface-viewpoint}
Theorem~\ref{thm:tholog-hypersurface} says that anomaly completion replaces \(B\) by a one-parameter degree-\(2\) hypersurface thickening.
The variable \(u\) is not an auxiliary decoration.
It is the ambient coordinate in which \([\Theta]\) cuts out the completed cohomology ring.
\end{remark}

\subsection{Genus-Clifford completion}
\label{subsec:tholog-genus-clifford}

Adding genuine surface data.
A genus-\(g\) surface carries \(2g\) basic one-cycles, and in the anomaly-completed setting their algebraic shadow is an odd Clifford algebra whose anticommutator is controlled by the secondary class \(u\).

Let
\[
H:=H_{1}(\Sigma_{g};k).
\]
Its intersection pairing is symplectic.
After parity shift, \(V:=\Pi H\) carries the corresponding symmetric bilinear form.
Choose a symplectic basis
\[
a_{1},b_{1},\dots,a_{g},b_{g},
\]
and write
\[
\alpha_{i}:=\Pi a_{i},\qquad \beta_{i}:=\Pi b_{i}.
\]

\begin{definition}[Genus-Clifford completion]
\label{def:tholog-genus-clifford}
The \emph{genus-\(g\) Clifford completion} of \(B_{\Theta}\) is the dg algebra
\[
\mathfrak{G}_{g}(B_{\Theta})
:=
\frac{
B_{\Theta}\langle \alpha_{1},\beta_{1},\dots,\alpha_{g},\beta_{g}\rangle
}{
\left\langle
\alpha_{i}\alpha_{j}+\alpha_{j}\alpha_{i},\;
\beta_{i}\beta_{j}+\beta_{j}\beta_{i},\;
\alpha_{i}\beta_{j}+\beta_{j}\alpha_{i}-\delta_{ij}u
\right\rangle
},
\]
with
\[
|\alpha_{i}|=|\beta_{i}|=1,\qquad d\alpha_{i}=d\beta_{i}=0.
\]
\end{definition}

\begin{remark}[Two regimes]
\label{rem:tholog-two-regimes}
Because the odd generators satisfy a Clifford relation with coefficient \(u\), the genus theory naturally splits into two regimes:
\begin{enumerate}
\item the \emph{curved} regime \(u\neq 0\), in which the Clifford algebra becomes Morita-trivial after localization;
\item the \emph{strict} regime \(u=0\), in which the handle algebra degenerates to an exterior algebra on \(H_{1}(\Sigma_{g})\).
\end{enumerate}
This dichotomy is one of the main structural consequences of the whole construction.
\end{remark}

\begin{theorem}[Universal property of genus completion; \ClaimStatusProvedHere]
\label{thm:tholog-universal-property-genus}
To give a left dg \(\mathfrak{G}_{g}(B_{\Theta})\)-module is equivalent to giving:
\begin{enumerate}
\item a left dg \(B_{\Theta}\)-module \(M\);
\item closed degree-\(1\) \(B_{\Theta}\)-linear endomorphisms
\[
A_{1},B_{1},\dots,A_{g},B_{g}\in \underline{\operatorname{End}}^{1}_{B_{\Theta}}(M)
\]
such that
\[
A_{i}A_{j}+A_{j}A_{i}=0,\qquad
B_{i}B_{j}+B_{j}B_{i}=0,\qquad
A_{i}B_{j}+B_{j}A_{i}=\delta_{ij}u_{M}.
\]
\end{enumerate}
\end{theorem}

\begin{proof}
This is immediate from generators and relations.
A module over \(\mathfrak{G}_{g}(B_{\Theta})\) is determined by the images of the generators
\(\alpha_{i},\beta_{i}\).
Preservation of the differential gives closedness of the corresponding endomorphisms, and preservation of the defining relations gives exactly the three displayed identities.
Conversely, any such family extends uniquely to a dg representation.
\end{proof}

\subsection{Morita triviality away from the secondary-anomaly cone}
\label{subsec:tholog-morita-triviality}

Let
\[
R:=B_{\Theta}[u^{-1}]
\]
be the localization of \(B_{\Theta}\) at the central element \(u\).

\begin{theorem}[One handle becomes a matrix factor after inverting \(u\); \ClaimStatusProvedHere]
\label{thm:tholog-one-handle-matrix}
There is an isomorphism of dg algebras
\[
\mathfrak{G}_{1}(R)\cong \underline{\operatorname{End}}_{R}
\big(R\oplus R\varepsilon\big),
\]
where \(|\varepsilon|=-1\).
\end{theorem}

\begin{proof}
Let
\[
S_{1}:=R\oplus R\varepsilon=R\otimes \Lambda^{\bullet}(k\varepsilon[-1]).
\]
Define degree-\(1\) endomorphisms of \(S_{1}\) by
\[
\rho(\alpha)=u\,(\varepsilon\wedge -),\qquad
\rho(\beta)=\iota_{\varepsilon}.
\]
Because
\[
(\varepsilon\wedge -)^{2}=0,\qquad \iota_{\varepsilon}^{2}=0,\qquad
(\varepsilon\wedge -)\iota_{\varepsilon}+\iota_{\varepsilon}(\varepsilon\wedge -)=\operatorname{id},
\]
we get
\[
\rho(\alpha)^{2}=0,\qquad
\rho(\beta)^{2}=0,\qquad
\rho(\alpha)\rho(\beta)+\rho(\beta)\rho(\alpha)=u\,\operatorname{id}.
\]
Hence \(\rho\) extends to a dg algebra map
\[
\mathfrak{G}_{1}(R)\longrightarrow \underline{\operatorname{End}}_{R}(S_{1}).
\]

As a left \(R\)-module, \(\mathfrak{G}_{1}(R)\) is free of rank \(4\) with basis
\[
1,\alpha,\beta,\alpha\beta.
\]
Likewise \(\underline{\operatorname{End}}_{R}(S_{1})\) is free of rank \(4\), with basis
\[
\operatorname{id},\quad \varepsilon\wedge -,\quad \iota_{\varepsilon},\quad
(\varepsilon\wedge -)\iota_{\varepsilon}.
\]
The image of \(\rho\) contains all four basis elements, because \(u\) is invertible and
\[
\varepsilon\wedge - = u^{-1}\rho(\alpha),\qquad
\iota_{\varepsilon}=\rho(\beta),\qquad
(\varepsilon\wedge -)\iota_{\varepsilon}=u^{-1}\rho(\alpha\beta).
\]
Thus \(\rho\) is an isomorphism.
\end{proof}

\begin{theorem}[Morita triviality of higher genus after inverting \(u\); \ClaimStatusProvedHere]
\label{thm:tholog-morita-triviality}
For every \(g\geq 0\),
\[
\mathfrak{G}_{g}(B_{\Theta})[u^{-1}]
\cong
\underline{\operatorname{End}}_{R}
\Big(R\otimes \Lambda^{\bullet}(k^{g}[-1])\Big),
\qquad
R=B_{\Theta}[u^{-1}].
\]
In particular,
\[
D\!\big(\mathfrak{G}_{g}(B_{\Theta})[u^{-1}]\text{-mod}\big)
\simeq
D\!\big(B_{\Theta}[u^{-1}]\text{-mod}\big).
\]
\end{theorem}

\begin{proof}
By construction, \(\mathfrak{G}_{g}(R)\) is the graded tensor product of \(g\) copies of \(\mathfrak{G}_{1}(R)\).
Likewise
\[
R\otimes \Lambda^{\bullet}(k^{g}[-1])
\cong
S_{1}^{\otimes g}.
\]
Taking graded tensor products of the isomorphism from Theorem~\ref{thm:tholog-one-handle-matrix} yields
\[
\mathfrak{G}_{g}(R)\cong
\underline{\operatorname{End}}_{R}(S_{1}^{\otimes g})
\cong
\underline{\operatorname{End}}_{R}
\Big(R\otimes \Lambda^{\bullet}(k^{g}[-1])\Big).
\]
Derived Morita equivalence is immediate.
\end{proof}

\begin{corollary}[Localization principle; \ClaimStatusProvedHere]
\label{cor:tholog-localization-principle}
All genuinely new genus-\(g\) information in \(\mathfrak{G}_{g}(B_{\Theta})\) is supported on the secondary-anomaly cone
\[
V(u)\subset \operatorname{Spec}H^{\bullet}(B_{\Theta}).
\]
Away from \(u=0\), the genus completion is derived Morita-trivial.
\end{corollary}

\subsection{Degeneration on the strict locus}
\label{subsec:tholog-degeneration-strict-locus}

\begin{theorem}[Exterior degeneration modulo \(u\); \ClaimStatusProvedHere]
\label{thm:tholog-exterior-degeneration}
There is an isomorphism of dg algebras
\[
\mathfrak{G}_{g}(B_{\Theta})/(u)
\cong
\big(B_{\Theta}/(u)\big)\otimes
\Lambda(\alpha_{1},\beta_{1},\dots,\alpha_{g},\beta_{g}),
\]
where all \(2g\) generators \(\alpha_{i},\beta_{i}\) have degree \(1\) and square to zero.
\end{theorem}

\begin{proof}
Reducing the defining relations modulo \(u\) gives
\[
\alpha_{i}\alpha_{j}+\alpha_{j}\alpha_{i}=0,\qquad
\beta_{i}\beta_{j}+\beta_{j}\beta_{i}=0,\qquad
\alpha_{i}\beta_{j}+\beta_{j}\alpha_{i}=0.
\]
This is exactly the exterior algebra relation.
\end{proof}

\begin{corollary}[Handle action in the strict sector; \ClaimStatusProvedHere]
\label{cor:tholog-handle-action-strict-sector}
If \(M\) is a left \(\mathfrak{G}_{g}(B_{\Theta})\)-module on which \(u\) acts by zero, then the handle operators define an exterior action
\[
\Lambda^{\bullet}\big(H_{1}(\Sigma_{g};k)\big)\curvearrowright M
\]
after parity shift.
\end{corollary}

\subsection{Handle operators on morphism complexes}
\label{subsec:tholog-handle-operators-on-morphisms}

\begin{theorem}[Handle action on morphism complexes; \ClaimStatusProvedHere]
\label{thm:tholog-handle-action-on-hom}
Let \(M,N\) be left \(\mathfrak{G}_{g}(B_{\Theta})\)-modules, with handle operators
\[
A_{i}^{M},B_{i}^{M}
\qquad\text{and}\qquad
A_{i}^{N},B_{i}^{N}.
\]
Define degree-\(1\) endomorphisms of \(\underline{\operatorname{Hom}}_{B_{\Theta}}(M,N)\) by
\[
\delta_{i}^{\alpha}(f)=A_{i}^{N}f-(-1)^{|f|}fA_{i}^{M},
\qquad
\delta_{i}^{\beta}(f)=B_{i}^{N}f-(-1)^{|f|}fB_{i}^{M}.
\]
Then
\[
(\delta_{i}^{\alpha})^{2}=0,\qquad
(\delta_{i}^{\beta})^{2}=0,
\]
\[
\delta_{i}^{\alpha}\delta_{j}^{\alpha}+\delta_{j}^{\alpha}\delta_{i}^{\alpha}=0,
\qquad
\delta_{i}^{\beta}\delta_{j}^{\beta}+\delta_{j}^{\beta}\delta_{i}^{\beta}=0,
\]
and
\[
\delta_{i}^{\alpha}\delta_{j}^{\beta}+\delta_{j}^{\beta}\delta_{i}^{\alpha}
=
\delta_{ij}\big(u_{N}\circ(-)-(-)\circ u_{M}\big).
\]
In particular, if \(u_{M}=u_{N}\), then the morphism complex carries a canonical action of the exterior algebra on \(2g\) generators.
\end{theorem}

\begin{proof}
Each identity is a direct graded expansion using the relations from Theorem~\ref{thm:tholog-universal-property-genus}.
For instance,
\[
(\delta_{i}^{\alpha})^{2}(f)=A_{i}^{N\,2}f-fA_{i}^{M\,2}=0,
\]
because \(A_{i}^{2}=0\).
Likewise,
\begin{align*}
\delta_{i}^{\alpha}\delta_{j}^{\beta}(f)+\delta_{j}^{\beta}\delta_{i}^{\alpha}(f)
&=
(A_{i}^{N}B_{j}^{N}+B_{j}^{N}A_{i}^{N})f
-
f(A_{i}^{M}B_{j}^{M}+B_{j}^{M}A_{i}^{M})
\\
&=
\delta_{ij}\big(u_{N}f-fu_{M}\big).
\end{align*}
The remaining identities are identical.
\end{proof}

\begin{remark}[The qualitative picture]
\label{rem:tholog-qualitative-picture}
Theorem~\ref{thm:tholog-handle-action-on-hom} is the algebraic form of a clean slogan:
\emph{equal secondary-anomaly sectors see genuine handle operators; different secondary-anomaly sectors see curvature instead of exterior symmetry}.
\end{remark}

\subsection{A strict holographic datum}
\label{subsec:tholog-strict-holographic-datum}

The theorem package above is purely algebraic.
The following definition packages it in the form suited to holographic applications.

\begin{definition}[Strict anomaly-completed holographic datum]
\label{def:tholog-strict-holographic-datum}
A \emph{strict anomaly-completed holographic datum} consists of:
\begin{enumerate}
\item a dg bialgebra \(B\), thought of as a bulk/defect/line algebra;
\item a primitive closed central class
\[
\Theta\in Z^{2}(B)\cap Z(B);
\]
\item a dg algebra \(A_{\partial}\), thought of as a boundary algebra;
\item a dg algebra map
\[
\beta:B\longrightarrow \underline{\operatorname{C}}^{\bullet}(A_{\partial},A_{\partial}),
\]
to the Hochschild cochain algebra of \(A_{\partial}\).
\end{enumerate}
A \emph{boundary neutralization} of this datum is a degree-\(1\) Hochschild cochain
\[
\gamma\in \underline{\operatorname{C}}^{1}(A_{\partial},A_{\partial})
\]
such that
\[
d\gamma=\beta(\Theta)
\]
and
\[
\gamma\,\beta(b)=(-1)^{|b|}\beta(b)\,\gamma
\qquad (b\in B \text{ homogeneous}).
\]
\end{definition}

\begin{remark}[Relation to the holographic modular Koszul datum]
\label{rem:tholog-hmkd-comparison}
\index{holographic modular Koszul datum!comparison with strict datum}
The holographic modular Koszul datum
$\mathcal H(T) = (\cA,\cA^!_{\mathrm{line}},\cC,r(z),\Theta_\cA,
\nabla^{\mathrm{hol}})$
of Volume~I
is the modular-operad packaging. The strict anomaly-completed
datum above is its dg-algebra shadow: the bulk/defect algebra~$B$
is a strict dg model for the line-side dg-shifted Yangian package
carried by $\cA^!_{\mathrm{line}}$, and the anomaly-completed line
algebra is the transgressed extension $B_{\Theta}$. The primitive class
$\Theta\in Z^2(B)$ is the image of the modular MC element
$\Theta_\cA$ under the collision residue
$\operatorname{Res}^{\mathrm{coll}}_{0,2}$
(from the collision-residue compatibility theorem of Volume~I),
the boundary algebra $A_\partial$ is the boundary chiral
algebra~$\cA$, and $\beta$ is the bulk-to-boundary map. The
boundary neutralization $\gamma$ with $d\gamma=\beta(\Theta)$
is the genus-$0$ incarnation of the shadow connection flatness:
at higher genus, $\gamma$ extends to
$\nabla^{\mathrm{hol}}_{g,n}=d-\operatorname{Sh}_{g,n}(\Theta_\cA)$
with $(\nabla^{\mathrm{hol}})^2=0$ from the full MC equation.
\end{remark}

\begin{proposition}[Universal boundary lift; \ClaimStatusProvedHere]
\label{prop:tholog-universal-boundary-lift}
To give a strict anomaly-completed holographic datum together with a boundary neutralization is equivalent to giving a dg algebra map
\[
\widetilde\beta:B_{\Theta}\longrightarrow \underline{\operatorname{C}}^{\bullet}(A_{\partial},A_{\partial}).
\]
\end{proposition}

\begin{proof}
This is exactly Proposition~\ref{prop:tholog-universal-property}, applied to the Hochschild cochain algebra of \(A_{\partial}\).
\end{proof}

\begin{remark}[Computed holographic data]
\label{rem:tholog-computed-examples}
\index{holographic modular Koszul datum!computed examples}
Volume~I computes the full holographic modular Koszul datum
$\mathcal{H}(T)$ for four families: Heisenberg, affine
Kac--Moody, Virasoro, and $\mathcal W_3$.
Under the comparison of
Remark~\ref{rem:tholog-hmkd-comparison}, the strict
anomaly-completed datum $(B,\Theta,A_\partial,\beta)$
specialises to:
\begin{center}
\renewcommand{\arraystretch}{1.15}
\begin{tabular}{llll}
& $B$ & $r^{\mathrm{coll}}(z)$ & $A_\partial$ \\
\hline
Heisenberg & $\Sym^{\mathrm{ch}}(V^*)$ & $k/z$ & $\cH_k$ \\
$\widehat{\mathfrak{sl}}_2{}_k$
 & $\widehat{\mathfrak{sl}}_2{}_{-k-4}$
 & $\Omega + k\delta/z$
 & $\widehat{\mathfrak{sl}}_2{}_k$ \\
$\mathrm{Vir}_c$
 & $\mathrm{Vir}_{26-c}$
 & $\frac{c/2}{z^3}+\frac{2T}{z}$
 & $\mathrm{Vir}_c$ \\
$\mathcal W_3{}_c$
 & $\mathcal W_3{}_{100-c}$
 & $r_c^{\mathrm{coll},\,W_3}(z)$
 & $\mathcal W_3{}_c$
\end{tabular}
\end{center}
\noindent
The column $r^{\mathrm{coll}}(z)$ is the bar-theoretic
collision residue (pole orders one lower than the Laplace kernel;
the Casimir $\Omega$ is a regular term at $z=0$).
The boundary neutralization $\gamma$ with
$d\gamma=\beta(\Theta)$ is the genus-$0$ shadow connection;
its higher-genus extension is
$\nabla^{\mathrm{hol}}_{g,n}=d-\mathrm{Sh}_{g,n}(\Theta_\cA)$.
For Virasoro at $c=26$, the intrinsic curvature is
$\kappa(\mathrm{Vir}_{26})=13$, but the effective curvature
$\kappa_{\mathrm{eff}}=(c-26)/2=0$ vanishes, so $\gamma$
extends to a flat connection at all genera.
For $\mathcal W_3$ at $c=100$, the same holds with
$\kappa_{\mathrm{eff}}(\mathcal W_3{}_{100})=0$.
\end{remark}

\begin{corollary}[Surface-completed holographic lift; \ClaimStatusProvedHere]
\label{cor:tholog-surface-completed-holographic-lift}
Suppose a strict anomaly-completed holographic datum is equipped with a boundary neutralization \(\gamma\), and suppose there exist closed degree-\(1\) Hochschild cochains
\[
A_{1},B_{1},\dots,A_{g},B_{g}
\]
satisfying
\[
A_{i}A_{j}+A_{j}A_{i}=0,\qquad
B_{i}B_{j}+B_{j}B_{i}=0,\qquad
A_{i}B_{j}+B_{j}A_{i}=\delta_{ij}\gamma^{2}.
\]
Then there is a dg algebra map
\[
\mathfrak{G}_{g}(B_{\Theta})
\longrightarrow
\underline{\operatorname{C}}^{\bullet}(A_{\partial},A_{\partial}).
\]
\end{corollary}

\begin{proof}
Apply Theorem~\ref{thm:tholog-universal-property-genus} to the dg algebra
\(\underline{\operatorname{C}}^{\bullet}(A_{\partial},A_{\partial})\),
viewed as a left module over itself.
\end{proof}

\subsection{Topological holography as an anomaly-completed Koszul-duality principle}
\label{subsec:tholog-koszul-duality-principle}

The strict algebra above suggests the following programme-level formulation.

\begin{principle}[Anomaly-completed topological holography]
\label{prin:tholog-anomaly-completed-topological-holography}
Let \(A_{\partial}\) be a boundary chiral or factorization algebra arising from a holomorphic-topological or twisted field theory.
Suppose its bulk/defect/line algebra is controlled by a Koszul-dual dg algebra \(B\simeq A_{\partial}^{!}\), and suppose there is a primitive closed central class
\[
\Theta_{g}\in Z^{2}(B)\cap Z(B)
\]
capturing the genus-\(g\) obstruction.
Then the correct strict genus-\(g\) bulk/defect/line algebra is not merely \(B\), but the anomaly-completed algebra
\[
B_{\Theta_{g}},
\]
and the correct surface-refined algebra is
\[
\mathfrak{G}_{g}(B_{\Theta_{g}}).
\]
\end{principle}

This principle is not claimed here as a theorem of physics.
What \emph{is} proved here is that \(B_{\Theta_{g}}\) and \(\mathfrak{G}_{g}(B_{\Theta_{g}})\) are the unique strict dg constructions with the universal properties one would want in such a theory.
They package:
\begin{itemize}
\item chain-level neutralizations of the genus class;
\item exact formulas for derived morphisms and relative tensor products;
\item a secondary anomaly controlling higher-genus support;
\item a genus-localization dichotomy: Morita-triviality away from \(u=0\), exterior degeneration on \(u=0\).
\end{itemize}

\subsection{A strict dg-shifted Yangian completion}
\label{subsec:tholog-strict-dg-shifted-yangian}

The line-operator literature motivating this appendix often packages OPE by a strict dg model of a shifted Yangian-type algebra.
The anomaly-completed form of that structure is immediate from the primitive lift theorem.

\begin{theorem}[Strict dg-shifted Yangian completion; \ClaimStatusProvedHere]
\label{thm:tholog-dg-shifted-yangian-completion}
Let \(Y\) be a strict dg model of a dg-shifted Yangian: a dg bialgebra equipped with translation automorphisms \(\tau_{z}\), twisted coproducts \(\Delta_{z}\), and a Maurer--Cartan \(r\)-matrix \(r(z)\).
Suppose \(\Theta\in Z^{2}(Y)\cap Z(Y)\) satisfies
\[
\tau_{z}(\Theta)=\Theta,\qquad
\Delta_{z}(\Theta)=\Theta\otimes 1+1\otimes \Theta,\qquad
\varepsilon(\Theta)=0.
\]
Then \(Y_{\Theta}\) carries a canonical strict dg-shifted Yangian structure extending that of \(Y\), determined by
\[
\tau_{z}(\eta)=\eta,\qquad
\Delta_{z}(\eta)=\eta\otimes 1+1\otimes \eta,\qquad
\varepsilon(\eta)=0,
\]
and with the same \(r\)-matrix \(r(z)\).
\end{theorem}

\begin{proof}
Translation invariance of \(\eta\) is compatible with the differential because
\[
d(\tau_{z}(\eta))=d(\eta)=\Theta=\tau_{z}(\Theta)=\tau_{z}(d\eta).
\]
Compatibility of the coproduct with the differential follows from the primitive condition:
\[
d\,\Delta_{z}(\eta)=\Theta\otimes 1+1\otimes \Theta=\Delta_{z}(\Theta)=\Delta_{z}(d\eta).
\]
The defining relation \(\eta y=(-1)^{|y|}y\eta\) with \(y\in Y\) is preserved exactly as in the proof of Theorem~\ref{thm:tholog-primitive-lift}.
Since \(r(z)\) already lies in \(Y\otimes Y\subset Y_{\Theta}\otimes Y_{\Theta}\) and the differential on \(Y\) is unchanged, the Maurer--Cartan and Yang--Baxter identities remain valid.
\end{proof}

\begin{remark}[Interpretation]
\label{rem:tholog-yangian-interpretation}
Theorem~\ref{thm:tholog-dg-shifted-yangian-completion} says that if the line side is modeled by modules for a strict dg-shifted Yangian before anomaly completion, then the correct genus-anomaly-corrected line algebra is \(Y_{\Theta}\), not \(Y\).
The \(r\)-matrix survives unchanged, while the module theory, Ext theory, and surface refinements are altered by the transgression generator \(\eta\) and its square \(u=\eta^{2}\).
\end{remark}

\subsection{Derived-center reconstruction and interval transgression}
\label{subsec:tholog-derived-center-reconstruction}

The surrounding literature repeatedly points to the derived center \(HH^{\bullet}(A_{\partial})\) as the natural recipient of bulk classes and to derived tensor products as the correct algebraic form of interval compactification.
The strict theorems above sharpen both themes.

\begin{proposition}[Boundary reconstruction of the anomaly; \ClaimStatusProvedHere]
\label{prop:tholog-boundary-reconstruction}
Suppose a strict anomaly-completed holographic datum admits a quasi-isomorphism
\[
\beta_{\mathrm{der}}:B\stackrel{\sim}{\longrightarrow} HH^{\bullet}(A_{\partial})
\]
onto the derived center of \(A_{\partial}\).
Then the anomaly class \(\Theta\) is completely determined by the boundary class
\[
\Theta_{\partial}:=\beta_{\mathrm{der}}(\Theta)\in HH^{2}(A_{\partial}),
\]
and for every \(A_{\partial}\)-module \(M\), the obstruction to neutralization is the image of \(\Theta_{\partial}\) in
\[
\operatorname{Ext}^{2}_{A_{\partial}}(M,M).
\]
\end{proposition}

\begin{proof}
The class \(\Theta\) is recovered from \(\Theta_{\partial}\) by the quasi-isomorphism \(\beta_{\mathrm{der}}\).
The action of the derived center on any module \(M\) sends \(\Theta_{\partial}\) to the obstruction class measuring whether the action of \(\Theta\) is null-homotopic on \(M\).
This is exactly Corollary~\ref{cor:tholog-obstruction-class} transported across the quasi-isomorphism.
\end{proof}

\begin{corollary}[Interval transgression formula; \ClaimStatusProvedHere]
\label{cor:tholog-interval-transgression-formula}
Let \(L,R\) be boundary conditions modelled by right and left modules over a strict bulk algebra \(B_{\Theta}\).
Then the derived-tensor core of the interval compactification is
\[
L\otimes^{L}_{B_{\Theta}}R
\simeq
\operatorname{Cone}\!\left(
L\otimes^{L}_{B}R[-1]
\xrightarrow{\;\theta_{L,R}\;}
L\otimes^{L}_{B}R
\right),
\]
with \(\theta_{L,R}\) as in Theorem~\ref{thm:tholog-derived-tensor}.
\end{corollary}

\begin{proof}
This is Theorem~\ref{thm:tholog-derived-tensor}.
\end{proof}

\subsection{Shadow calculus: semisimple, Casimir, and scalar families}
\label{subsec:tholog-shadow-calculus}

The strict algebraic formalism above is universal.
To recover concrete module computations one passes to a shadow in which the derived endomorphism complex is computable.
The following general lemmas extract the three patterns that repeatedly appear in the surrounding examples: semisimple vanishing, Casimir control, and scalar universality.

\begin{proposition}[Semisimple vanishing criterion; \ClaimStatusProvedHere]
\label{prop:tholog-semisimple-vanishing}
Let \(M\) be a left dg \(B\)-module.
Assume that there exists a finite-dimensional semisimple Lie algebra \(\mathfrak{g}\) and a finite-dimensional \(\mathfrak{g}\)-module \(E\) such that
\[
R\!\operatorname{End}_{B}(M)\simeq C^{\bullet}(\mathfrak{g},E)
\]
in the derived category of dg algebras.
Then any central degree-\(2\) anomaly class acts trivially on \(M\):
\[
[\Theta_{M}]=0\in \operatorname{Ext}^{2}_{B}(M,M).
\]
\end{proposition}

\begin{proof}
Under the hypothesis, the obstruction class lands in
\[
H^{2}(\mathfrak{g},E).
\]
Whitehead's second lemma gives
\[
H^{2}(\mathfrak{g},E)=0
\]
for semisimple \(\mathfrak{g}\) and finite-dimensional \(E\).
Hence the obstruction vanishes.
\end{proof}

\begin{remark}[Evaluation-family shadow]
\label{rem:tholog-evaluation-family-shadow}
Whenever a family of modules (for example an evaluation-family model of a Yangian sector) has derived self-endomorphisms computed by the Chevalley--Eilenberg complex of a finite-dimensional semisimple Lie algebra, Proposition~\ref{prop:tholog-semisimple-vanishing} applies immediately.
This recovers the expected vanishing phenomenon in a precise hypothesis-driven form.
\end{remark}

\begin{proposition}[Casimir shadow lemma; \ClaimStatusProvedHere]
\label{prop:tholog-casimir-shadow}
Let \(C\) be an algebra equipped with a central element \(\Omega\), and let \(S\) be a simple left \(C\)-module on which \(\Omega\) acts by the scalar \(\chi_{S}(\Omega)\).
Suppose that a shadow functor sends the anomaly class \(\Theta_{g}\) to the scalar multiple
\[
\lambda_{g}\Omega
\]
in the center of \(C\).
Then the induced shadow of the anomaly on \(S\) is multiplication by
\[
\lambda_{g}\chi_{S}(\Omega).
\]
\end{proposition}

\begin{proof}
A central element acts on a simple module by a scalar.
If the shadow of \(\Theta_{g}\) is \(\lambda_{g}\Omega\), then its action is exactly the displayed scalar multiple.
\end{proof}

\begin{remark}[Affine Zhu shadow]
\label{rem:tholog-affine-zhu-shadow}
In the standard affine setting, the relevant shadow algebra is Zhu's algebra \(U(\mathfrak{g})\), and \(\Omega\) is the quadratic Casimir.
If the genus anomaly descends to
\[
\frac{k+h^{\vee}}{2h^{\vee}}\lambda_{g}\Omega,
\]
then on a simple lowest-weight shadow \(L_{\mathfrak{g}}(\bar\lambda)\) it acts by
\[
\frac{k+h^{\vee}}{2h^{\vee}}\,c_{2}(\bar\lambda)\,\lambda_{g}.
\]
This Casimir law's precise realization depends on the ambient affine model.
\end{remark}

\begin{proposition}[Scalar shadow lemma; \ClaimStatusProvedHere]
\label{prop:tholog-scalar-shadow}
Suppose that in a given shadow category the anomaly class \(\Theta_{g}\) acts by a scalar
\[
\kappa\,\lambda_{g}
\]
independently of the chosen simple object in a family.
Then every simple object in that family has the same shadow anomaly:
\[
[\Theta_{g}]=\kappa\,\lambda_{g} \qquad \]
\end{proposition}

\begin{proof}
This is immediate from the hypothesis.
\end{proof}

\begin{remark}[Fock-family shadow]
\label{rem:tholog-fock-family-shadow}
The Heisenberg/Fock pattern is exactly of the form treated by Proposition~\ref{prop:tholog-scalar-shadow}: the genus anomaly is universal across the family.
This is the sharp contrast with semisimple evaluation sectors, where the anomaly vanishes, and with affine Casimir shadows, where it depends on the lowest-weight Casimir eigenvalue.
\end{remark}

\subsection{The finished first-principles picture}
\label{subsec:tholog-finished-picture}

We may now summarize the whole construction in one paragraph.

Start with a dg algebra \(B\) and a closed central degree-\(2\) class \(\Theta\).
Adjoin a skew odd transgression generator \(\eta\) with \(d\eta=\Theta\); this produces the anomaly-completed algebra \(B_{\Theta}\).
The universal property of \(B_{\Theta}\) identifies its modules with \(B\)-modules equipped with chain-level neutralizations of \(\Theta\).
The corrected one-sided two-term resolution computes both
\(R\!\operatorname{Hom}_{B_{\Theta}}\) and derived relative tensor products by explicit fiber and cone formulas.
The square \(u=\eta^{2}\) is a new closed central class, the \emph{secondary anomaly}; it endows the space of neutralizations with a natural quadratic refinement and controls the curvature of the transgression operator on morphism complexes.
Passing from \(B_{\Theta}\) to \(\mathfrak{G}_{g}(B_{\Theta})\) adjoins \(2g\) handle operators with Clifford anticommutator \(u\).
After inverting \(u\), this higher-genus completion becomes Morita-trivial.
Modulo \(u\), it degenerates to an exterior extension by \(H_{1}(\Sigma_{g})\).
Thus the genuinely new genus-\(g\) information is concentrated exactly on the secondary-anomaly cone \(u=0\).

This is, in a strict algebraic sense, the anomaly-completed topological-holographic package toward which the surrounding chapters point.

\subsection{Frontier conjectures and repository-facing statements}
\label{subsec:tholog-frontier-conjectures}

The strict M-level algebra proved in this appendix suggests several precise H-level conjectures.
We record them in a form suitable for the research repository.

\begin{conjecture}[Koszul-dual bulk completion; \ClaimStatusConjectured]
\label{conj:tholog-koszul-dual-bulk-completion}
Let \(A_{\partial}\) be a boundary chiral algebra arising from a 3d
holomorphic-topological theory
(Theorem~\ref{thm:physics-bridge}), and
let \(A_{\partial}^{!}\) be its open-colour Koszul dual. On the
chirally Koszul locus, the line category is modeled by
\(A_{\partial}^{!}\)-modules by
Theorem~\textup{\ref{thm:lines_as_modules}}.
Then:
\begin{enumerate}[label=\textup{(\roman*)}]
\item Volume~I's genus-\(g\) curvature class
 \(\kappa(\cA) \cdot \omega_g \in H^2_{\mathrm{cyc}}(\cA,\cA)\)
 maps, via the Hochschild bridge
 (Theorem~\textup{\ref{thm:hochschild-bridge-higher-genus}}),
 to a canonical class
 \(\Theta_g \in Z^2(A_{\partial}^{!}) \cap Z(A_{\partial}^{!})\).
\item \(\Theta_g\) is central because \(\kappa(\cA)\) is scalar
 (Volume~I, Theorem~D, the modular characteristic):
 centrality is the shadow of the period-corrected differential
 \(\Dg{g}\) being a coderivation
 (Volume~I, Remark~\ref{V1-rem:sc-higher-genus}\textup{(iii)}),
 which the Hochschild bridge transports to a central cocycle
 in \(A_\partial^!\).
\item \(\Theta_g\) is primitive in \(B_\Theta^{(g)}(\cA)\)
 after passing to the period-corrected differential
 \(\Dg{g}\): the fibrewise differential \(\dfib\) is
 \emph{not} a coderivation of \(\Delta\) at higher genus
 (Volume~I, Remark~\ref{V1-rem:sc-higher-genus}\textup{(ii)}),
 but the Fay trisecant period correction restores both
 flatness and the coderivation property simultaneously, after
 which the Hodge curvature class
 \(\kappa(\cA)\cdot\omega_g\) pushes forward to a primitive
 element (Proposition~\textup{\ref{prop:R-canonical-vol2}(iii)},
 in its corrected form).
 The Hochschild bridge transports primitivity.
\item The genus-\(g\) bulk algebra is therefore the transgression
 algebra \((A_{\partial}^{!})_{\Theta_g}\) of
 Definition~\textup{\ref{def:tholog-transgression-algebra}}, with
 all the strict algebraic structure proved in
 Sections~\textup{\ref{subsec:tholog-transgression-algebra}--\ref{subsec:tholog-handle-operators-on-morphisms}}.
\end{enumerate}
\end{conjecture}

\begin{remark}[Heuristic justification]
The genus-\(g\) curvature \(\kappa(\cA) \cdot \omega_g\) is an
element of \(H^2_{\mathrm{cyc}}(\cA, \cA)\) by Volume~I,
Theorem~D. The Hochschild bridge
(Theorem~\ref{thm:hochschild-bridge-higher-genus}) provides a
filtered quasi-isomorphism
\(\mathrm{C}^\bullet_{\mathrm{ch,top}}(\cA) \simeq
\ChirHoch^\bullet(\cA)\) at genus~\(g\). The cyclic cohomology
class \(\kappa(\cA) \cdot \omega_g\) maps under the open-colour
Koszul duality functor to a class
\(\Theta_g \in H^2(A_\partial^!)\): this uses the isomorphism
\(H^2_{\mathrm{cyc}}(A_\partial, A_\partial) \cong
\mathrm{Ext}^2_{A_\partial^!\text{-}\mathrm{bimod}}(A_\partial^!, A_\partial^!)\)
on the Koszul locus (Volume~I, Theorem~H, at generic level; the critical
level $k = -h^\vee$ is excluded because $\dim \ChirHoch^0$ can be infinite
there, see Vol~I Remark~\ref*{rem:critical-level-lie-vs-chirhoch}).

Centrality: \(\kappa(\cA)\) is scalar (acts as multiplication by a
constant on every bar component), so \(\Theta_g\) commutes with
all elements of \(A_\partial^!\). Primitivity: by
Proposition~\ref{prop:curved-R-factorization}\textup{(iii)} the
fibrewise differential \(\dfib\) is \emph{not} a coderivation of
\(\Delta\) at higher genus, but the period-corrected
differential \(\Dg{g}\) of Volume~I,
Remark~\ref{V1-rem:sc-higher-genus}\textup{(iii)}, is both flat
and a coderivation; under \(\Dg{g}\), the curvature class
\(\kappa(\cA)\cdot\omega_g\) acts on cogenerators by the scalar
\(\kappa(\cA)\) and is therefore primitive in
\(B_\Theta^{(g)}(\cA)\). The Hochschild bridge preserves both
centrality and primitivity because it is a filtered
quasi-isomorphism respecting the bialgebra structure.

With \(\Theta_g\) established as a primitive central \(2\)-cocycle,
the transgression algebra \((A_{\partial}^{!})_{\Theta_g}\) is
defined by Definition~\ref{def:tholog-transgression-algebra}, and
all the M-level structure theorems
(Sections~\ref{subsec:tholog-transgression-algebra}--\ref{subsec:tholog-handle-operators-on-morphisms})
apply.
\end{remark}

\begin{conjecture}[Surface-completed line-side algebra; \ClaimStatusConjectured]
\label{conj:tholog-surface-completed-line-algebra}
In the setting of Conjecture~\ref{conj:tholog-koszul-dual-bulk-completion},
the genus-\(g\) refinement of the line-side algebra is
\[
\mathfrak{G}_{g}\!\big((A_{\partial}^{!})_{\Theta_{g}}\big).
\]
In addition:
\begin{enumerate}[label=\textup{(\roman*)}]
\item The full higher-genus line category is supported on the
 secondary-anomaly cone \(u_{g} = \eta^2 = 0\).
\item Away from \(u_g = 0\), the genus completion is derived
 Morita-trivial
 (Theorem~\textup{\ref{thm:tholog-morita-triviality}}).
\end{enumerate}
\end{conjecture}

\begin{remark}[Heuristic justification]
By Conjecture~\ref{conj:tholog-koszul-dual-bulk-completion},
\(\Theta_g\) is a primitive central \(2\)-cocycle in \(A_\partial^!\).
The genus-Clifford completion
\(\mathfrak{G}_g((A_\partial^!)_{\Theta_g})\)
is defined by attaching \(g\) pairs of odd generators
\((\eta_\alpha, \bar\eta_\alpha)\) satisfying the Clifford
relation \(\eta_\alpha \bar\eta_\alpha +
\bar\eta_\alpha \eta_\alpha = u_g\), where
\(u_g = \eta^2\) is the secondary anomaly
(Proposition~\ref{prop:tholog-secondary-central-class}).

On the chirally Koszul locus, the line-category comparison
\(\mathcal{C}_{\mathrm{line}} \simeq \cA^!_{\mathrm{line}}\text{-}\mathbf{mod}\)
(Theorem~\ref{thm:lines_as_modules}) lifts to
the genus-completed setting: at genus~\(g\), the genus-completed
line side is modeled by modules for
\(\mathfrak{G}_g((A_\partial^!)_{\Theta_g})\).

Part~(i): by the localization principle
(Corollary~\ref{cor:tholog-localization-principle}), all genuinely
new genus-\(g\) information is supported on \(u_g = 0\).
Part~(ii): Theorem~\ref{thm:tholog-morita-triviality} gives
\(\mathfrak{G}_g((A_\partial^!)_{\Theta_g})[u_g^{-1}] \simeq
(A_\partial^!)_{\Theta_g}[u_g^{-1}]\) (Morita equivalence via
matrix factorization after inverting~\(u_g\)).
\end{remark}

\begin{conjecture}[Boundary reconstruction; \ClaimStatusConjectured]
\label{conj:tholog-boundary-reconstruction}
Assume the bulk--boundary identification
\(A_{\partial}^{!} \simeq HH^{\bullet}(A_{\partial})\)
holds (Theorem~\textup{\ref{thm:hochschild-bridge-higher-genus}}
on the Koszul locus). Then:
\begin{enumerate}[label=\textup{(\roman*)}]
\item The genus anomaly \(\Theta_{g}\) is reconstructed from a
 canonical boundary Hochschild class
 \(\Theta_{g,\partial} \in HH^{2}(A_{\partial})\):
 the Hochschild bridge transports the genus-\(g\) curvature
 \(\kappa(\cA) \cdot \omega_g\) to
 \(\Theta_{g,\partial}\), and the bulk--boundary equivalence
 maps \(\Theta_{g,\partial}\) to~\(\Theta_g\).
\item The surface-completed bulk algebra
 \(\mathfrak{G}_g((A_\partial^!)_{\Theta_g})\) is equivalent to
 the universal genus-completed endomorphism algebra of the boundary
 theory: it controls exactly those deformations of the boundary
 algebra that extend to genus-\(g\) surfaces.
\end{enumerate}
\end{conjecture}

\begin{remark}[Heuristic justification]
Part~(i): the chain of identifications is
\[
\kappa(\cA) \cdot \omega_g
\;\xmapsto{\text{Hochschild bridge}}\;
\Theta_{g,\partial} \in HH^2(A_\partial)
\;\xmapsto{A_\partial^! \simeq HH^\bullet(A_\partial)}\;
\Theta_g \in Z^2(A_\partial^!) \cap Z(A_\partial^!).
\]
The first arrow is
Theorem~\ref{thm:hochschild-bridge-higher-genus}(iii) (curvature
maps to cyclic obstruction). The second arrow is the
Koszul-locus identification of \(A_\partial^!\) with
\(HH^\bullet(A_\partial)\) (this is exactly
Proposition~\ref{prop:tholog-boundary-reconstruction},
transported across the quasi-isomorphism). Together, the
genus anomaly is determined entirely by boundary data.

Part~(ii): the transgression algebra
\((A_\partial^!)_{\Theta_g}\) controls deformations of
\(A_\partial^!\)-modules that neutralize the anomaly class~\(\Theta_g\)
(Corollary~\ref{cor:tholog-modules-neutralizations}). Under the
bulk--boundary equivalence, this is the same as deformations of
the boundary theory that extend from the disk (genus~$0$) to a
genus-\(g\) surface, precisely because $\Theta_g$ measures the
obstruction to this extension. The genus-Clifford completion
\(\mathfrak{G}_g\) adds the handle operators, giving the
universal endomorphism algebra that controls all genera
simultaneously.
\end{remark}

\begin{remark}[The conjectural anomaly-holographic machine]
\label{rem:tholog-complete-machine}
Conjectures~\ref{conj:tholog-koszul-dual-bulk-completion}--\ref{conj:tholog-boundary-reconstruction}
form the H-level completion target for the anomaly-holographic machine of this chapter.
The strict M-level algebra
(Sections~\ref{subsec:tholog-transgression-algebra}--\ref{subsec:tholog-handle-operators-on-morphisms})
applies to every boundary chiral algebra arising from a 3d HT
theory (Theorem~\ref{thm:physics-bridge}), and the conjectures
above describe the additional bridge data needed to identify that
strict model with the physical bulk/line package:
\begin{enumerate}[label=\textup{(\alph*)}]
\item the homotopy-Koszulity of \(\SCchtop\)
 (Theorem~\ref{thm:homotopy-Koszul}) gives the open-colour
 Koszul dual \(A_\partial^!\);
\item the Hochschild bridge
 (Theorem~\ref{thm:hochschild-bridge-higher-genus}) gives
 \(\Theta_g\);
\item the M-level structure theorems give the full transgression,
 genus-Clifford, and localization package.
\end{enumerate}
The boundary fixes the formal-local bulk model at all genera.
In physical realizations, identifying that strict model with the
bulk package remains the H-level bridge problem: the proposed
holographic dictionary is presented by the transgression algebra,
as the chiral derived center pair is presented by the bar-cobar package.
\end{remark}

\begin{remark}[Repository integration]
\label{rem:tholog-repository-integration}
For repository purposes, the clean separation is:
\begin{itemize}
\item Sections~\ref{subsec:tholog-transgression-algebra}--\ref{subsec:tholog-handle-operators-on-morphisms} are theorem-level strict algebra proved here.
\item Sections~\ref{subsec:tholog-strict-holographic-datum}--\ref{subsec:tholog-derived-center-reconstruction} are strict packaging results and exact consequences.
\item Section~\ref{subsec:tholog-shadow-calculus} records representation-theoretic shadows in a hypothesis-driven form.
\item Conjectures~\ref{conj:tholog-koszul-dual-bulk-completion}--\ref{conj:tholog-boundary-reconstruction} record the H-level target: the M-level algebra is proved here, while extending it to all HT theories requires the additional bridge identifications stated above.
\end{itemize}
\end{remark}

\begin{remark}[Extension to $\mathfrak{sl}_4$: cross-orbit Koszul duality; \ClaimStatusProvedHere]
\label{rem:tholog-sl4-extension}
\index{anomaly completion!sl4@$\mathfrak{sl}_4$}
The $\mathfrak{sl}_3$ pattern proved in the Bershadsky--Polyakov computation
(Computation~\ref{comp:bp-anomaly-completion} in the frontier chapter)
extends to $\mathfrak{sl}_4$ with qualitatively new features.
The nilpotent poset of $\mathfrak{sl}_4$ has five orbits:
$(1^4)\subset(2,1^2)\subset(2,2)\subset(3,1)\subset(4)$.
The three intermediate orbits produce three distinct anomaly-completion
patterns:
\begin{enumerate}[label=\textup{(\roman*)}]
\item \emph{Subregular $(3,1)$}: five strong generators of weights
 $1,2,2,2,3$; cross-orbit Koszul dual to the minimal $(2,1,1)$.
\item \emph{Rectangular $(2,2)$}: seven strong generators of
 weights $1^3,2^4$; same-orbit Koszul dual ($(2,2)^t=(2,2)$),
 the first non-hook, non-principal orbit in the programme.
\item \emph{Minimal $(2,1,1)$}: nine strong generators of
 weights $1^4,(3/2)^4,2$; cross-orbit Koszul dual to the
 subregular $(3,1)$.
\end{enumerate}
The pair $(3,1)\leftrightarrow(2,1,1)$ is the first genuinely
cross-orbit Koszul duality in the anomaly-completed programme:
the two $\mathcal W$-algebras have different rank, different
generator spectra, and different ghost systems ($3$ vs $4$ pairs).
The transgression algebra $B_\Theta$ bridges both simultaneously.

The curvature for all orbits satisfies
$\kappa(\mathcal W_\lambda^k)=\frac{15(k+4)}{8}
-\kappa_{\mathrm{ghost}}^{(\lambda)}(k)$,
confirming the $(k+h^\vee)$ proportionality law at leading order.
A \emph{JM parity criterion} governs level-dependence:
orbits with even-integer Jacobson--Morozov grading ($(3,1)$, $(2,2)$)
have level-independent $\kappa_{\mathrm{ghost}}$ and exact
proportionality $\kappa\propto(k+4)$; the minimal orbit $(2,1,1)$
has odd-integer JM grades and level-dependent
$\kappa_{\mathrm{ghost}}(k)$.
The full computation is given in
Conjecture~\ref{comp:sl4-hook-anomaly-completion}.
\end{remark}

\subsection{Coda}
\label{subsec:tholog-coda}

Topological holography, at the strict algebraic level accessible here, is not an isomorphism between bulk and boundary operator algebras.
It is a \emph{Koszul-dual anomaly-completed transgression theory}.
The classical/commutative/coisson bulk class \(\Theta\) is not merely an obstruction; it is the source of a universal odd transgression \(\eta\), a central secondary charge \(u=\eta^{2}\), and a whole higher-genus Clifford/exterior machine.
The higher-genus structure is therefore not uniformly new:
away from the secondary-anomaly cone it collapses Morita-theoretically,
while on the cone it produces a genuine exterior handle algebra.

%=======================================================================
% Anomaly-completed Yangian identification
%=======================================================================

\subsection{The anomaly-completed Yangian identification}
\label{subsec:anomaly-completed-yangian-identification}

The transgression data $(\Theta,\eta,u)$ determines a deformation of
the line category $\cA^!_{\mathrm{line}}$ that we identify, under the
affine monodromy identification of
Theorem~\ref{thm:affine-monodromy-identification} (in
\texttt{log\_ht\_monodromy\_core.tex}), with a rigidly
\emph{anomaly-completed Yangian}.

\begin{theorem}[Anomaly-completed Yangian identification]
\label{thm:anomaly-completed-yangian}
\index{Yangian!anomaly-completed}
\index{transgression!Yangian pair}
Let $\cA = V_k(\fg)$ be the affine Kac--Moody vertex algebra at
non-critical level $k\neq -h^\vee$. Let $(\Theta,\eta,u)$ be the
anomaly-completed transgression data of
Part~\ref{part:anomaly-completed}, with $u = \eta^2$ the secondary
central charge. Then:
\begin{enumerate}[label=\textup{(\roman*)}]
\item The line category $\cA^!_{\mathrm{line}}$ carries a canonical
$u$-deformed Yangian structure $Y_\hbar(\fg;u)$, specializing to
the standard Yangian $Y_\hbar(\fg)$ at $u=0$.
\item The pair $(Y_\hbar(\fg;u),\Theta)$ is the physical
Yangian-plus-shadow datum: $Y_\hbar(\fg;u)$ is the line/boundary
category, $\Theta\in\mathrm{MC}(\mathfrak{g}^{\mathrm{mod}}_\cA)$ is
the shadow obstruction tower controlling its genus expansion, and
$\eta$ is the odd transgression connecting them.
\item At the self-dual anomaly point $u = u^*$, the deformation
$Y_\hbar(\fg;u)$ degenerates to the Clifford completion of
Theorem~\ref{V1-thm:g9-clifford-structure}; this is the unique
point where the celestial shadow class is $\mathbf C$
(Theorem~\ref{thm:class-C-quartic-bms}).
\end{enumerate}
\end{theorem}

\begin{proof}
(i) The $u$-deformation is constructed as follows: the secondary
central extension by $u=\eta^2$ modifies the coproduct of
$Y_\hbar(\fg)$ by an additive $u$-cocycle, yielding a Hopf
algebra deformation $Y_\hbar(\fg;u)$. This is the unique
deformation compatible with the transgression structure
(uniqueness: the secondary central charge is unique up to
scale by the Koszul-dual anomaly theory of
\S~\ref{subsec:tholog-coda}). (ii) Identification of $\Theta$ with
the shadow obstruction tower is
Theorem~\ref{thm:novikov-completion-theorem} in
\texttt{ym\_instanton\_screening.tex}; identification of $\eta$ as
the connecting transgression is the anomaly-completion theorem of
Part~\ref{part:anomaly-completed}. (iii) Clifford degeneration at
$u=u^*$: at the self-dual point, the $u$-deformed coproduct
acquires a Clifford-type quartic relation, which is the algebraic
signature of class~$\mathbf C$.
\end{proof}

\begin{remark}[Physical content: Y and $\Theta$ are one pair]
\label{rem:anomaly-yangian-one-pair}
The physical upshot of Theorem~\ref{thm:anomaly-completed-yangian}
is that the Yangian (line algebra) and the shadow obstruction
tower (genus expansion) are not independent: they are two sides
of a single $u$-deformed structure, bridged by the transgression
$\eta$. This resolves the apparent tension between the algebraic
line category (finite-dimensional at each level) and the
gravitational genus expansion (infinite), by identifying them as
complementary faces of the anomaly-completed data. In particular,
the Costello--Paquette $\cN=4$ celestial algebra is the $u = u^*$
specialization (class~$\mathbf C$) of this construction.
\end{remark>
